L'usage de l'algorithme de Fisher-Yates permet de piocher $m$ arcs aléatoires de parmi ceux à éventuellement supprimer. La structure Unir \& Trouver représente sur l'ensemble des arêtes (donc des murs) la relation "séparer la même pièce". Ainsi, pour vérifier que la suppression d'un mur ne crée pas de cycle, il suffit de vérifier que les deux pièces séparées par ce mur ne sont pas déjà dans la même classe d'équivalence. En effet, si l'on supprime un mur séparant deux pièces déjà dans la même classe d'équivalence, alors il existe déjà un chemin entre elles. En supprimant le mur, on crée un cycle. Il n'y a donc grâce à la structure Unir \& Trouver aucun risque de créer des boucles, ni des cavités\footnote{une cavité désigne une pièce contenant au moins un carré $2 \times 2$.}.\par
Pour ce qui est des salles isolées, on le doit au fait que toutes les arêtes, au nombre de $m = 2n(n-1)$ sont testées. Ainsi, si dans le parcours une pièce était isolée, alors le parcours des arêtes concernées entraînerait la suppression d'au moins un mur, car deux composantes connexes seront fusionnées avec succès. Par conséquent, il n'y a pas de risque d'avoir des salles isolées.\par
On obtient à la fin un arbre couvrant du graphe des pièces du labyrinthe. Ainsi, il existe un unique chemin entre deux pièces du labyrinthe. En particulier, l'entrée et la sortie sont reliées par un chemin unique.
